\newpage
\section{Déroulement du projet}
\label{sec:equipe}

Dans cette section, nous décrivons comment le projet a été réalisé en équipe : la répartition des tâches, la synchronisation du travail en membres de l'équipe, etc.

\subsection{Réalisation du projet par étapes}

Le projet a été conduit selon un découpage en étapes successives, en s'inspirant d'un cycle de développement incrémental.

\paragraph{Étape 1 — Analyse et cahier des charges} Étude du sujet, rédaction du cahier des charges, identification des grandes entités (maison, pièces, meubles, maître, besoins, routine) et premières maquettes papier de l'interface.

\paragraph{Étape 2 — Squelette du moteur} Mise en place du paquetage \texttt{config}, de la classe \texttt{Map}, des classes \texttt{Room} et de leurs filles, ainsi que des classes de meubles. Création du \texttt{GameBuilder} pour assembler la carte et les meubles.

\paragraph{Étape 3 — État du maître et horloge} Implémentation des quatre besoins (\texttt{Energy}, \texttt{Hunger}, \texttt{Thirst}, \texttt{Hygiene}), de \texttt{MasterState} (humeur, besoin critique) et de \texttt{SimulationClock}.

\paragraph{Étape 4 — Routine et imprévus} Conception des classes \texttt{RoutineStep}, \texttt{RoutinePlan} et \texttt{RoutineManager}, avec deux modèles (semaine et week-end). Ajout du \texttt{ImprevuManager} et des cinq imprévus (\textit{Grasse matinée}, \textit{Visite surprise}, \textit{Coupure de courant}, \textit{Journée malade}, \textit{Journée pressée}).

\paragraph{Étape 5 — Décision et navigation} Mise en place du pipeline de décision (\texttt{DecisionMaker} avec trois \texttt{DecisionStrategy}) et du \texttt{MasterNavigationService}.

\paragraph{Étape 6 — Sélecteurs d'humeur et services} Ajout des sélecteurs (\texttt{MoodFoodSelector}, \texttt{MoodShowSelector}) et du \texttt{FurnitureInteractionService} basé sur le patron \textit{Visitor}.

\paragraph{Étape 7 — Interface graphique} Conception du \texttt{LauncherGUI}, du \texttt{MainGUI} et de tous les panneaux (\texttt{GameDisplay}, \texttt{StatsPanel}, \texttt{FurnitureStatusPanel}, etc.).

\paragraph{Étape 8 — Tests et qualité} Rédaction de la suite de tests \texttt{DomotiqueTestSuite}, qui regroupe une vingtaine de classes de tests \texttt{JUnit} couvrant les besoins, l'horloge, la routine, les imprévus, les sélecteurs, les meubles, la navigation et l'intégration via le \texttt{MobileElementManager}.

\paragraph{Étape 9 — Rédaction du rapport} Mise en forme du rapport en \LaTeX, intégration des schémas et relectures.

\subsection{Répartition des tâches, communication et intégration}

\section{Répartition des tâches, communication et intégration}

\subsection{Répartition des tâches}

Le projet a été divisé en plusieurs grands domaines fonctionnels, chacun confié à un ou plusieurs membres de l'équipe selon leurs affinités et compétences.

\begin{itemize}
    \item \textbf{Gestion des besoins et système de décision}~: développement des états du personnage (énergie, faim, soif, hygiène), mise en place du système de routine et implémentation du mécanisme de priorisation entre besoins critiques et routine.
    \item \textbf{Interface graphique et rendu visuel}~: conception de la fenêtre principale, réalisation des panneaux d'affichage (statistiques, état des meubles, notifications), intégration des sprites et des animations du personnage.
    \item \textbf{Navigation et interactions}~: gestion des pièces et des meubles, implémentation des déplacements du maître et de la logique d'interaction avec les objets de la maison.
\end{itemize}

Chaque membre a pu contribuer à l'ensemble du projet tout en étant référent sur son domaine principal. Cette répartition a permis une progression parallèle et efficace.

\subsection{Communication}

La communication au sein du groupe a été régulière et fluide. Plusieurs outils ont été utilisés pour faciliter les échanges et le suivi du projet~:

\begin{itemize}
    \item Un serveur \texttt{Discord} dédié pour les discussions quotidiennes, le partage rapide de documentation et la résolution de problèmes en temps réel.
    \item Des réunions hebdomadaires en présentiel pour faire le point sur l'avancement, valider les choix d'architecture et répartir les tâches à venir.
    \item Un dépôt \texttt{Git} centralisé pour versionner le code et suivre l'historique des modifications.
\end{itemize}

Cette organisation a maintenu une bonne coordination et a permis de respecter les délais impartis.

\subsection{Intégration}

L'intégration des différents modules a été réalisée de manière progressive selon une approche itérative~:

\begin{enumerate}
    \item Développement individuel sur des branches dédiées.
    \item Tests unitaires et validation locale des fonctionnalités.
    \item Fusion régulière sur la branche principale lors des sessions de travail collectif.
\end{enumerate}

Cette méthode a facilité la détection précoce des incohérences et a garanti la cohérence globale de l'application. L'ensemble des composants a ainsi pu être intégré avec succès pour livrer une simulation fonctionnelle répondant aux exigences initiales.

